\section{Experiment Execution}
\subsection{Subjects Preparation}
In order to not just look at workloads from two different microservice systems, but also at different levels of complexity and behaviour patterns within these microservice systems, two usage scenarios are defined for each system.
The two usage scenarios for \textit{Sock Shop} and \textit{Train Ticket} are comparable across systems. The first scenario for \textit{Sock Shop} involves opening the catalogue and randomly selecting an item for preview. In the second scenario, the user then logs in, opens the catalogue and adds a random number of items to cart to create an order with. For \textit{Train Ticket}, a user adds users to the database, logs in and then searches for tickets in the first scenario. In the second scenario, the user again adds users to the database and logs in, but this time they search for a train ticket and book it. Across both scenarios and applications, users randomly wait 1-3 seconds between carrying out any action. Furthermore, the random numbers used throughout these scenarios are always picked from a range of 3 to 9. To script these scenarios, Locust.io\footnote{\url{https://locust.io}} is used.

\subsection{Setup}
\subsubsection{Collecting Ground Truth Measures}
To collect the ground truth measures, \textit{Sock Shop} and \textit{Train Ticket} are deployed on two servers provided by the \textit{Green Lab}\footnote{\url{https://s2group.cs.vu.nl/pages/greenlab}}. These servers each run Ubuntu 22.04.2 LTS on Intel Xeon Silver 4208 CPU @2.10 GHz with 16 cores and 384 GB of RAM. 
On these servers, Scaphandre\footnote{\url{https://github.com/hubblo-org/scaphandre}} is used to collect energy consumption measurement values in $\mu W$ of the services individually and the system overall every minute. Scaphandre does so by combining the per-socket RAPL with the DRAM RAPL power measures. To collect metrics regarding resource usage such as CPU and memory usage, cAdvisor\footnote{\url{https://github.com/google/cadvisor}} is used.
In order to make sure it does not interfere with the measurement, the collection of those measures is done on a third server running on an Intel Xeon Processor E3-1230 v3 @ 3.40 GHz with 8 vCPUs, 32 GB of RAM and 1 TB of HD, also operating on Ubuntu 22.04 LTS. These three servers are connected to a switch (see \autoref{fig:servers}) and share a bandwidth of approximately 1 Gbps and a latency of approximately 0.25 ms. For the collection of measures from the other two servers, Prometheus\footnote{\url{https://prometheus.io}} is run on this server. The systems on the servers are deployed using Docker version 27.3.1.

\begin{figure}
    \centering
    \includegraphics[width=0.6\linewidth]{vu-cs-research-thesis/figures/serverSetup.png}
    \caption{Setup of the servers used to collect ground truth measures.}
    % \Description{Setup of the servers used to collect ground truth measures.}
    \label{fig:servers}
\end{figure}

\subsubsection{Simulation Executions}
Using Python version 3.10.4, swf and parquet traces are produced from the obtained resource mesures for each treatment consisting of one of the microservice applications, scenarios, user loads and repetitions. From the Locust logs, some high-level communication traces are produced. Given that distributed tracing is not supported for the microservice applications, only requests to and from the applications are logged. 

Execution of the simulations is done on an additional workstation server containing Intel Xeon Silver 4208 CPU @ 2.10 GHz with 16 cores and 384 GB of RAM. This was to run each simulator in a best case scenario where they were not limited by resources. The simulators were run on Ubuntu 22.04.2 LTS with Python 3.10.4, Java (OpenJDK) 21.0.7, and g++ 11.4.0.

\subsection{Execution Phase 1: Ground Truth Collection}
For collecting the ground truth, every trial is run 30 times in total. Each experimental run begins with a three minute warm-up period to make sure stable idle cpu utilization levels. A treatment is then randomly selected, ensuring that the order of trials does not bias the results. The trial's user scenario is then run via locust for 5 minutes, after which data is collected with Prometheus. Finally, a cooldown period of 2 minutes is initiated to ensure the system returns to idle levels. The warm-up and cooldown periods ensure that trials do not influence each other and are as accurate and reliable as possible. Prometheus is used to collect power consumption metrics from Scaphandre at 5-second intervals as well as service-level metrics from cAdvisor such as CPU utilization and memory usage. These values are queried from Prometheus using PromQL\footnote{\url{https://prometheus.io/docs/prometheus/latest/querying/basics}}.

\subsubsection{Trace Construction}
The values queried from Prometheus are saved and later used to generate traces. In order to reduce the large number of traces ($2 \cdot 2 \cdot 2 \cdot 30 = 240$), for each treatment, the four repetitions with the total energy consumption closest to the median are selected. The resulting selection of traces is then simulated with the configurations described in \autoref{tab:config-vars}, resulting in $32 \cdot (4 + 2 + 4 + 5 + 1) = 512$ simulation instances that are each repeated 30 times in order to account for any possible non-deterministic behaviour from the simulators. In particular, CPU and memory utilization are used to generate \texttt{swf} and \texttt{parquet} traces, as well as defining the described workloads in code. For the \texttt{swf} traces, each 5-second measurement of a service is defined as a job in order to preserve the degree of granularity of the CPU and memory utilization measures.

For the communication traces, for each trace, the mean response size is calculated and saved. The communication trace itself consists of the number of requests per second time step.

\subsubsection{Ground Truth Construction}
The power consumption samples collected from Scaphandre using Prometheus is used to build the ground truth dataset against which the simulated energy consumption values are compared. As the simulators produce energy consumption values on a host-level, we are using the power samples that describe the energy usage of the whole system. In order to do so, the obtained power consumption values are transformed from $\mu W$ to Joules consumed in that time step as described in \autoref{eq:energy_trapezoid}. The power consumption values are transformed from $\mu W$ to $W$. Then, the trapezoid rule is applied by computing the mean of the power values at the beginning ($p_{i-1}$) and end ($p_i$) of each time step, and multiplying this by the duration of the time step in seconds to estimate the energy consumed during that interval and adding this value to the calculated energy consumed during the previous intervals. Each service is given a number and jobs are allocated to the service by giving this number both as its group ID and executable number.

% \begin{equation} \label{eq:energy}
%     \text{Energy}(t_\text{begin}, t_\text{end}) = \frac{p_\text{begin} + p_\text{end}}{2} \cdot (t_\text{end} - t_\text{begin}) \cdot 10^{-6}
% \end{equation}

\begin{equation} \label{eq:energy_trapezoid}
\text{Energy}(t_0, t_N) 
= \int_{t_0}^{t_n} p(t)\, dt
\approx 10^{-6} \sum_{i=0}^{n} \frac{p_{i-1} + p_i}{2} (t_i - t_{i-1})
\end{equation}

\begin{table*}[ht]
\centering
\caption{Mean relative error for the selected simulators per power model. Bold values are a simulator's lowest MRE.}
\label{tab:powermodelOverview}
\begin{tabular}{l|rrrrrr|r}
 & \textit{\textbf{CloudSim}} & \textit{\textbf{NetworkCloudSim}} & \textit{\textbf{OpenDC}} & \textit{\textbf{DISSECT-CF}} & \multicolumn{2}{l|}{\textit{\textbf{SimGrid}}} & \multirow{2}{*}{\textbf{Mean}} \\
 & code & \begin{tabular}[l]{@{}l@{}}code with\\ communication\end{tabular} & parquet & swf & code & \begin{tabular}[l]{@{}l@{}}code with\\ communication\end{tabular} &  \\ \hline
linear          & 0.74          & 0.74          & 0.41          & 0.71          & \textbf{0.42} & \textbf{0.42} & 0.57 \\
square          & 0.63          & 0.63          & 0.41          & -             & -             & -             & 0.56 \\
cubic           & \textbf{0.54} & \textbf{0.54} & 0.41          & -             & -             & -             & 0.50 \\
square root     & 0.80          & 0.81          & 0.37          & -             & -             & -             & 0.66 \\
constant        & -             & -             & \textbf{0.14} & \textbf{0.44} & -             & -             & 0.29 \\ \hline
Mean            & 0.68          & 0.68          & 0.35          & 0.57          & 0.42          & 0.42          & 
\end{tabular}
\end{table*}

\subsection{Execution Phase 2: Simulations Executions} \label{sec:simTestingExecution}
In order to execute the simulations on the selected simulators, the treatments described in Table \ref{tab:config-vars} are each run 30 times to increase reliability and deal with any intrinsic non-determinism of the simulators. The simulated energy usage for each run is written to file along with the runtime of each simulation.

\subsubsection{System Properties Collection}
In order to be able to define the systems used to collect measures for simulation, some properties have to be collected. As a measure of CPU performance, the Instruction Per Cycle (IPC) or variations of it are needed for all simulators except \textit{OpenDC}. This metric is not just dependent on the system but also the applications running. In order to get an approximation of this value, the IPC of the servers running the applications and monitoring tools was calculated using the Linux \verb|perf stat| command over 1 minute and repeated 30 times for both \textit{Sock Shop} and \textit{Train Ticket}. The mean of these 30 IPC values was then taken to be used to configure the cloud simulators. For configuration of the simulators' power models, power consumption measures at three different CPU utilization levels have to be collected: \textit{off}, \textit{idle} and \textit{max}. The off level refers to the powerdraw when it is completely off. As we are collecting the energy usage from the server itself and not its external powerdraw for this study, this value is set to 0. For the idle level, we take the mean of host power values for the time steps of the ground truth dataset for which the sum of CPU usage over all monitored containers is 0\%. For the max level, we do the same but for a CPU usage of 100\% or more, which occurs without additional interference.

\subsubsection{CloudSim}
Workloads are defined in code for \textit{CloudSim}. In order to do so, a \verb|PowerContainer|, \verb|PowerContainerVM| and \verb|Cloudlet| are created for each service and a custom \verb|UtilizationModel| is defined to provide the CPU utilization per time step for the given service to the \verb|Cloudlet| from the trace. While \textit{CloudSim} does provide a way to load workloads as \texttt{swf} files, it only takes the id, runtime and number of processors fields into account, meaning that e.g. task submission in time is not possible. Because of this, we decided not to simulate workloads as swf files for \textit{CloudSim}.

\subsubsection{NetworkCloudSim} 
The manner in which workloads defined in code for \textit{NetworkCloudSim} is similar to that for \textit{CloudSim}. For each service, a \verb|PowerContainer| and \verb|PowerContainerVM| are created, as well as a \verb|NetworkCloudlet|. The custom \verb|UtilizationModel| defined for \textit{CloudSim} is reused. For the communication, two extra \verb|NetworkCloudlets| are created representing the requests and responses. For each request from the trace, a request and response with the mean response size are added to these two \verb|NetworkCloudlets|. In order to be able to run the \verb|NetworkCloudlets| in combination with energy modelling, the custom \verb|PowerNetworkHost|, \verb|PowerSwitch| and \verb|PowerContainerNetworkDatacenter| classes are defined. These classes only make the combination of the power and network modules and do not add extra capabilities or logic.

\subsubsection{DISSECT-CF}
Workloads are given to \textit{DISSECT-CF} as \texttt{swf} trace. To load the \texttt{swf}, the provided \verb|FileBasedTraceProducerFactory| is used. As the energy measures that the simulated values will be compared to are only from the CPU, the powerstates of the other resources such as storage and network are set to 0 and only the consumed energy of the hosts are registered.

\subsubsection{OpenDC}
Workloads are given to \textit{OpenDC} as two \texttt{parquet} files. The services are defined in the \textit{tasks} file, while the CPU usage metrics of the services over time are given in the \textit{fragments} file. With \textit{OpenDC}, it is possible to execute multiple workload simulators in batch during one experiment, in this case of 8 workloads. In order to obtain the simulation time of each workload, the simulation time of the experiment is divided by the number of workloads simulated.

\subsubsection{SimGrid}
For SimGrid, the workloads are given in code, both with and without communication. This is done using the S4U interface and the \textit{Host Energy} plugin. For each service, an S4U agent is defined with a function that reads the CPU consumption values for that service. From the CPU consumption value, the actual runtime of the CPU for that time step is calculated and the service is set to sleep for the remainder of the time step. To add communication, an extra agent is added on the original host which responds to requests with the mean response size. To send these requests an additional host was created with an agent to generate them. The communication between these agents was modelled using the synchronous communication interface provided by SimGrid. Only the original host's energy consumption is recorded.


